\documentclass[UTF8,a4paper,11pt,openany]{ctexbook}
\usepackage{../common/statlearnbook}
\SLSetBookMetadata{第 5 册}{统计学习方法讲义 v2.6.0}{采样方法、主题模型与图排序}
\begin{document}
\setcounter{page}{729}
\setcounter{chapter}{31}
\setcounter{figure}{1}
\pagestyle{headings}
\chaptermark{蒙特卡罗方法与直接采样}
\sectionmark{蒙特卡罗估计}
若还要使用有限方差误差公式，则必须再有
\[
  \int_{\{p>0\}}\frac{h(x)^2}{p(x)}\,\mathrm dx<\infty.
\]
所以“选任意密度就能计算任意积分”不是正确命题；密度必须覆盖被积函数的非零区域，且比值必须满足相应可积条件。阅读图\ref{fig:V5-C02-mc-integral}时应具体回答：被积函数非零而 $p=0$ 的区域会在哪一步丢失，以及二阶矩条件为何比一阶积分改写多出一个 $1/p$ 因子。图只表示这两个逻辑接口，不提供随机复算坐标。
% v2.3.1 figure UID: FIG-P573-01
% Proposed post-v2.3.1 polish for FIG-P573-01: protect key-label size and stroke hierarchy.
\tikzset{slfig-FIG-P573-01/.style={
  every path/.append style={line cap=round,line join=round},
  every node/.style={font=\fontsize{9.5pt}{11.4pt}\selectfont}
}}
\pgfplotsset{slfig-FIG-P573-01-axis/.style={
  tick label style={font=\fontsize{8.6pt}{10.3pt}\selectfont},
  label style={font=\fontsize{9.5pt}{11.4pt}\selectfont},
  axis line style={line width=.62pt},tick style={line width=.52pt},clip=false
}}
\begin{figure}[htbp]
\centering
\begin{tikzpicture}[slfig-FIG-P573-01]
\begin{axis}[slfig-FIG-P573-01-axis,width=.72\linewidth,height=6.0cm,xmin=0,xmax=1,ymin=0,ymax=1.08,
  name=mcaxis,
  axis lines=left,xlabel={$u$},ylabel={$h(u)=\exp(-u^2/2)$},
  xtick={0,.2,.4,.6,.8,1},ytick={0,.2,.4,.6,.8,1},
  clip=false]
  \addplot[name path=curve,draw=SLBlue,line width=1.05pt,domain=0:1,samples=301]{exp(-x^2/2)};
  \addplot[name path=axiszero,draw=none,domain=0:1]{0};
  \addplot[fill=SLBlue,fill opacity=.08,draw=none] fill between[of=curve and axiszero];
  \addplot[ycomb,draw=SLTeal,line width=.70pt,mark=*,mark size=2.2pt]
    coordinates {(.10,.9950125)(.40,.9231163)(.70,.7827045)(.80,.7261490)};
  \addplot[draw=SLGold,line width=.86pt] coordinates {(0,.8567456)(1,.8567456)};
  \addplot[draw=SLTeal,line width=.58pt,dash pattern=on 5pt off 2.4pt]
    coordinates {(0,.8556244)(1,.8556244)};
  \node[anchor=south west,font=\fontsize{9.5pt}{11.4pt}\selectfont,
    text=SLGold,fill=white,inner sep=1.4pt]
    at (axis cs:.45,.99) {样本均值 $\widehat\mu=0.8567$};
  \node[anchor=south west,font=\fontsize{9.5pt}{11.4pt}\selectfont,
    text=SLTeal,fill=white,inner sep=1.4pt]
    at (axis cs:.62,.89) {真实积分 $\mu=0.8556$（长虚线）};
\end{axis}
  \node[draw=SLRuleGray,line width=.55pt,rounded corners=2pt,fill=white,
    font=\fontsize{9.5pt}{11.4pt}\selectfont,inner xsep=3.4mm,inner ysep=2.2mm,
    anchor=north]
    at ([yshift=-10mm]mcaxis.south)
    {$\displaystyle\widehat\mu_N=\frac1N\sum_{i=1}^{N}h(u_i)$};
\end{tikzpicture}
\caption{蒙特卡罗积分把曲线下的面积改写为分布下的期望；竖线表示固定的四个均匀样本，虚线表示该积分的参考值}\label{fig:V5-C02-mc-integral}
\end{figure}

\textbf{阅读边界。}四个样本坐标是冻结的说明性点，不声称来自随机种子。
\end{document}

